\section{Representation-theoretic tools}
\label{sec:prelim}

We fix the tensor conventions used by both the coordinate projections
and the staircase input. The two transfer mechanisms below are dual
polytabloid detection of a constituent and ordinary induction on disjoint
position sets.

All representations are finite-dimensional over $\C$. A partition is
padded with zeros when needed; $\lambda^t$ denotes its conjugate.
For partitions of the same integer, $\lambda\unrhd\mu$ means
$\sum_{i\le k}\lambda_i\ge\sum_{i\le k}\mu_i$ for every $k$.
Conjugation reverses dominance. Every symmetric-group representation
is semisimple, every Specht module is self-dual, and
$S^\lambda\otimes\sgn\cong S^{\lambda^t}$.
We use the standard Specht construction, Young's rule, branching,
Schur--Weyl duality, and the Littlewood--Richardson and Pieri rules
\cite{James,FultonHarris,Macdonald}.
Their precise consequences needed below are recorded here.

\subsection{Alternating blocks and dual tests}

We realize Specht modules by column alternation and detect constituents
by contracting with dual tensors. This viewpoint also underlies
Ikenmeyer's triangular construction \cite[Sections~3--5]{Ikenmeyer}.

Let $B$ be a finite set of positions, and let $E$ have ordered basis
$e_1,\ldots,e_D$. For an ordered block $A=(b_1,\ldots,b_k)$, $k\le D$,
put
\[
 v_A=\sum_{\pi\in S_k}\sgn(\pi)
       e_{\pi(1)}\otimes\cdots\otimes e_{\pi(k)},
\]
with the displayed factors placed at $b_1,\ldots,b_k$.
For a set partition $\mathcal A$ of $B$, put
$v_{\mathcal A}=\bigotimes_{A\in\mathcal A}v_A$.
Changing the order of one block changes only the overall sign.
The position action of $S_B$ is the usual permutation of tensor factors;
all regroupings of factors are ordinary tensor identifications and add
no signs.

If the block sizes are the column lengths of $\theta$, then
\begin{equation}\label{eq:block-specht}
 \C[S_B]v_{\mathcal A}\cong S^\theta.
\end{equation}
Indeed, arrange the blocks as tableau columns, and send a row tabloid
to the word assigning letter $i$ to its $i$th row. Column alternation
sends its polytabloid to $v_{\mathcal A}$, which is nonzero because the
initial row word has coefficient one. The same construction applies
in a dual alphabet; see \cite[Definitions~4.1 and~4.3,
Theorem~4.12]{James} for the Specht construction and ordinary
irreducibility.

\begin{lemma}[Dual-polytabloid test]\label{lem:dual-test}
Let $z\in E^{\otimes B}$ and $\theta\vdash |B|$ with at most $D$ rows.
Then $S^\theta$ occurs in $\C[S_B]z$ if and only if $z$ has a nonzero
contraction with a dual block tensor whose blocks are the columns of
$\theta$, for some placement of the blocks and some common ordered
basis of $E^*$.
If such a contraction exists for a fixed placement, it is nonzero for
a generic ordered basis. Finitely many separately nonzero contractions
can be made nonzero simultaneously with a common generic basis.
\end{lemma}

\begin{proof}
For any fixed basis and placement the orbit of the dual block tensor
is a copy of $S^\theta$, by \eqref{eq:block-specht}.
Restricting the natural tensor pairing gives an equivariant map from
$\C[S_B]z$ to the dual of this copy of $S^\theta$.
A nonzero contraction makes this map nonzero, hence proves the required
constituent.
Conversely, Schur--Weyl duality
\cite[Appendix~A, (5.5$'$), (5.6), and (8.2)]{Macdonald}
identifies the $\theta$-isotypic
component with $S^\theta\otimes\mathbb S_\theta(E)$.
The symmetric-group translates and common basis changes of one
nonzero dual polytabloid span the corresponding full dual isotypic
component: each of the two factors is irreducible for its own group.
They therefore detect any nonzero $\theta$-component of $z$.
For a fixed placement, contraction is a polynomial in the entries of
the common change-of-basis matrix. A nonzero polynomial remains
nonzero on a nonempty Zariski-open subset of $\GL_D(\C)$.
The product of finitely many such polynomials is nonzero, proving
the final statement.
\end{proof}

\subsection{Support under induction}

Throughout, combining constituents on disjoint position sets means
\emph{ordinary induction} from the product of their symmetric groups,
not a Kronecker product. We use $\boxtimes$ for an outer tensor product.
Let $c_{\alpha,\beta}^{\gamma}$ denote the Littlewood--Richardson
coefficient, equivalently the multiplicity of $S^\gamma$ in
$\Ind(S^\alpha\boxtimes S^\beta)$.
This is the induction form of the LR rule
\cite[Rule~16.4]{James}; its polynomial $\GL$ counterpart is
\cite[Appendix~A, (6.1)]{Macdonald}.
Only support, rather than equality of multiplicities, is required.

\begin{lemma}[Young support]\label{lem:young-support}
For a finite list $A$ of positive integers of sum $n$, the module
$\Ind_{\prod_{a\in A}S_a}^{S_n}\mathbf1$ contains exactly the
$S^\nu$ whose column heights $h_1\ge h_2\ge\cdots$ satisfy
\begin{equation}\label{eq:young-prefix}
 \sum_{j=1}^k h_j\le P_A(k):=\sum_{a\in A}\min(k,a)
 \qquad(k\ge1).
\end{equation}
The constituent indexed by the decreasing rearrangement of $A$
has multiplicity one. A component containing every shape of total
$t$ with at most $j$ rows contains the support of the induced-trivial
module corresponding to a balanced division of $t$ into $j$ parts.
\end{lemma}

\begin{proof}
Young's rule \cite[Rule~14.1]{James} says that the support consists of the partitions
dominating the decreasing rearrangement of $A$, with diagonal
multiplicity one. Transposing dominance gives
\eqref{eq:young-prefix}. The balanced partition of $t$ into $j$ parts
is dominated by every partition of $t$ with at most $j$ rows;
conversely, any partition dominating it has at most $j$ rows.
Zero parts in a balanced division can be omitted.
\end{proof}

\begin{lemma}[Horizontal addition]\label{lem:lr-addition}
If $c_{\alpha,\beta}^{\gamma}>0$ and
$c_{\alpha',\beta'}^{\gamma'}>0$, then
\[
 c_{\alpha+\alpha',\,\beta+\beta'}^{\gamma+\gamma'}>0,
\]
where addition is rowwise. In particular, suppose a target has column
heights $h_i=a_i+b_i$, with $a_i,b_i\ge0$. Form $\alpha$ from the
column-height list $(a_i)$ and $\beta$ from $(b_i)$, omitting zeros
and sorting each list. Then $c_{\alpha,\beta}^{\gamma}>0$ for the
target $\gamma$.
\end{lemma}

\begin{proof}
The first assertion is the Littlewood--Richardson semigroup property
\cite[Theorem~1]{Zelevinsky}, applied to the two triples of partitions.
For the second assertion, the exterior-power case of Pieri gives
$c_{(1^{a_i}),(1^{b_i})}^{(1^{h_i})}>0$ for each $i$.
Add these containments horizontally. A rowwise sum of single-column
partitions is exactly the diagram specified by their column-height
multiset, so sorting the individual lists changes nothing.
\end{proof}

We also use two other immediate forms of the LR rule. First,
$c_{\alpha,(t)}^\gamma>0$ precisely when $\gamma/\alpha$ is a
horizontal $t$-strip, meaning at most one new cell per column.
This is Pieri's rule \cite[Chapter~I, (5.16)--(5.17)]{Macdonald}.
Second, if $\alpha$ fits in a rectangle, its rotated complement
$\beta$ satisfies $c_{\alpha,\beta}^{\text{rectangle}}=1$.
For a rectangle $(w^j)$, the complement has parts
$\beta_i=w-\alpha_{j+1-i}$, and the Schur-character formula gives
$\mathbb S_\beta(\C^j)\cong\det^w\otimes
\mathbb S_\alpha(\C^j)^*$
\cite[Chapter~I, (3.1); Appendix~A, (8.2) and (8.4)]{Macdonald}.
Since $\mathbb S_{(w^j)}(\C^j)=\det^w$, the LR tensor-product
rule and Schur's lemma give the asserted multiplicity one.
The zero-area case uses the empty partition, and the stated induction
conventions include empty partitions as well.

\subsection{The cyclic staircase input}

For $m\ge1$, set
\[
 B_m=\{(i,j)\in\Z_{\ge1}^2:i+j\le m+1\},\qquad
 N_m=\frac{m(m+1)}2,\qquad \rho_m=(m,m-1,\ldots,1).
\]
Let $R_m$ and $C_m$ be the actual row and column set partitions of
$B_m$. Order each row by increasing $j$ and each column by increasing
$i$. In each tensor layer use the ordered alphabet $e_1,\ldots,e_m$.
With the unnormalized initial-letter alternations defined above, set
\begin{equation}\label{eq:cyclic-generator}
 w_m=v_{R_m}\otimes v_{C_m},\qquad
 W_m=\C[S_{B_m}]w_m\subseteq S^{\rho_m}\otimes S^{\rho_m}.
\end{equation}
The action in \eqref{eq:cyclic-generator} is diagonal.

\begin{theorem}[Cyclic staircase support, {\cite[Theorem~3.1]{SaxlCyclic}}]
\label{thm:cyclic-input}
For every $m\ge1$ and every $\alpha\vdash N_m$, the irreducible
$S^\alpha$ occurs in the particular cyclic module $W_m$ of
\eqref{eq:cyclic-generator}.
\end{theorem}

The complete proof is given in Section~6 of the companion
\cite{SaxlCyclic}, with the same ordered row and column blocks as here.
The bibliography records the separate preprint's stable article identifier.
It is important that the theorem concerns the specified row/column
generator, not merely the ambient staircase square.
Common relabelings of the two ordered alphabets and changes of
block orders preserve the assertion.
